\section{Further consequences}\label{sec:consequences}

The main theorem gives both numerical information about blocks and a
functorial identity. We first prove the local--global bounds stated in
\Cref{cor:block-consequences}, then record the more detailed numerical
classification and the categorical reformulation. These are consequences
of the numerical theorem through the cited results; none is an input
to the proof above.

\subsection{Local--global and principal-block bounds}

The comparison with a defect-group normalizer is local to an individual
block. Summing it over maximal-defect blocks gives the comparison with
the Sylow normalizer, while the two principal-block bounds relate the
number of simple modules to the number of classes of $p$-elements.

We use the notation of \Cref{cor:block-consequences}; in particular,
$P\in\operatorname{Syl}_p(G)$.

\begin{proof}[Proof of \Cref{cor:block-consequences}]
The local inequality and its abelian-defect equality are Alperin's Consequences~1--2
\cite{Alperin1987}, as restated in \cite[Section~3, p.~334]{HST2024};
their numerical weight-conjecture hypotheses follow from \Cref{thm:main}.
Put $N=N_G(P)$. Since $P$ is a normal Sylow $p$-subgroup of $N$,
it is a defect group of every block of $N$.
Brauer's first main theorem therefore matches all $p$-blocks of $N$
with the maximal-defect $p$-blocks of $G$; this indexing is also used in
\cite[proof of Lemma~2.14]{FMR2025}. Writing these two block sets as
$\operatorname{Bl}_p(N)$ and $\operatorname{Bl}_{p,\max}(G)$, respectively,
the local inequality gives
\[
 |\IBr_p(G)|
 \geq\sum_{B'\in\operatorname{Bl}_{p,\max}(G)}l(B')
 \geq\sum_{b'\in\operatorname{Bl}_p(N)}l(b')
 =|\IBr_p(N)|.
\]
Finally, \Cref{thm:main} supplies the only conjectural input in
Hung--Sambale--Tiep's Propositions~3.1--3.2 \cite{HST2024}, which give
the two principal-block bounds under their respective hypotheses.
\end{proof}

\subsection{Character counts and defect data}

The next consequence treats blocks whose ordinary-character count
exceeds their Brauer-character count by one. Besides giving bounds,
it specifies all possible pairs
of defect-group order and Brauer-character count, and asserts that
every listed pair is attained.

\begin{corollary}[Character-count and defect data]\label{cor:character-data}
Let $B$ be a $p$-block of a finite group with defect group $D$ of order
$p^d$, and write $k(B)=|\Irr(B)|$. If $k(B)-l(B)=1$, then the possible
numerical data $(p^d,l(B))$ are exactly the following:
\begin{enumerate}[label=\textup{(\roman*)}]
\item There are positive integers $n,m$ such that $d=nm$, every prime divisor of $n$
      divides $p^m-1$, and $4\mid n$ implies $4\mid p^m-1$, with
      \[
       l(B)=\sum_{e\mid n}\frac{p^{em}-1}{ne}J_2(n/e),\qquad
       J_2(t)=t^2\prod_{\substack{q\mid t\\q\text{ prime}}}(1-q^{-2}).
      \]
\item The exceptional pairs are
      \[
       \begin{aligned}
       (p^d,l(B))\in\{&(5^2,7),(7^2,8),(11^2,9),(11^2,35),\\
                       &(23^2,88),(29^2,63),(59^2,261)\}.
       \end{aligned}
      \]
\end{enumerate}
All listed data occur, and in every case $k(B)=l(B)+1$.
\end{corollary}

\begin{proof}
Hung, Sambale, and Tiep proved this numerical classification assuming
the numerical weight conjecture for $B$ \cite[Theorem~3.3]{HST2024}.
\Cref{thm:main} supplies that hypothesis. Their examples establishing
that every listed value occurs were already unconditional.
\end{proof}

This classifies only character-count and defect numerical data, not blocks
up to isomorphism, Morita equivalence, or derived equivalence.

\subsection{Functorial formulation}

The following formulation expresses the numerical theorem as an identity
of classes of functors. Its alternating sum runs over all strict chains;
it is not the normal-chain transform used earlier to prove the theorem.

The numerical theorem gives the functorial formulation of
Boltje, Bouc, and Y\i lmaz \cite[Conjecture~3.3]{BBY2026}.
Let $k$ and $F$ be algebraically closed fields of characteristics $p$ and
zero, respectively. We recall the objects used in this optional
formulation. A finite-dimensional $(kH,kG)$-bimodule is viewed as a
$k[H\times G]$-module by $(h,g)m=hmg^{-1}$. It is a
$p$-permutation module if its restriction to a Sylow $p$-subgroup
has a basis permuted by that subgroup. Such a $p$-permutation bimodule
is diagonal if it is also
projective as a left $kH$-module and as a right $kG$-module.
The notation $T^\Delta(H,G)$ denotes their split Grothendieck group,
with relations $[M\oplus N]=[M]+[N]$.

The category of group--central-idempotent pairs has objects $(H,c)$,
including $c=0$. A morphism from $(H,c)$ to $(H',c')$ is an
$F$-linear combination of diagonal $p$-permutation bimodule classes
selected by $c'$ on the left and $c$ on the right. Composition is
induced by tensor product over the intermediate group algebra.
In the notation of \cite[Section~2]{BBY2026}, the category of
$F$-linear functors from this category to $F$-vector spaces is
$\mathcal F\mathcal B\ell_{F,k}$.

This functor category is semisimple in characteristic zero
\cite[Section~2.5]{BBY2026}. By Yoneda, the endomorphism space of
the representable $FT^\Delta(-,H)$ is the finite-dimensional
$F$-space $FT^\Delta(H,H)$. Semisimplicity therefore forces the
representable to have finite length: each simple summand contributes an
independent projection. Its direct summands also have finite length.
Here $K_0(\mathcal F\mathcal B\ell_{F,k})$ denotes the Grothendieck
group of the finite-length subcategory, the free abelian group on
simple-functor classes. It is in this sense that classes record the
simple multiplicities used in \cite[Sections~2.5 and~3]{BBY2026}.

The bimodule $kHc$ acts by precomposition on the representable
functor $FT^\Delta(-,H)$. Its image is the direct summand
$FT^\Delta_{H,c}=FT^\Delta(-,H)\circ kHc$ associated to $(H,c)$,
using the equivalent group and pair descriptions in
\cite[Sections~2.3--2.4]{BBY2026}. In particular $c=0$ gives the
zero functor. Write $[[H,c]]_F$ for its class in
$K_0(\mathcal F\mathcal B\ell_{F,k})$.

Write $S_{1,1,F}$ for the simple
functor indexed by the trivial pair and its trivial $F$-module.
Its multiplicity in $FT^\Delta_{H,c}$ is $l(H,c)$: evaluation at
the trivial group gives the space spanned by indecomposable projective
$kHc$-modules \cite[proof of Theorem~3.4]{BBY2026}. Thus the
following identity records all simple-functor multiplicities, with
the simple-module count as its $S_{1,1,F}$ component.

\begin{corollary}[Functorial Alperin weight identity]\label{cor:functorial-awc}
Let $b$ be a block idempotent of $kG$, where $G$ is finite. Let
$\mathcal S_p(G)$ consist of all strict chains
$\sigma=(1=P_0<P_1<\cdots<P_n)$ of $p$-subgroups, including $(1)$.
Put
\[
 |\sigma|=n,\qquad
 G_\sigma=\bigcap_{i=0}^nN_G(P_i),\qquad
 b_\sigma=\operatorname{Br}_{P_n}(b).
\]
Since $C_G(P_n)\leq G_\sigma$ and the Brauer image is invariant
under $G_\sigma$, the element $b_\sigma$ is a central idempotent
of $kG_\sigma$, possibly zero or a sum of blocks. Thus the pair
$(G_\sigma,b_\sigma)$ is an object of the category just described.
Then, summing over representatives of the $G$-conjugacy classes of chains,
\[
 \sum_{\sigma\in[G\backslash\mathcal S_p(G)]}
 (-1)^{|\sigma|}[[G_\sigma,b_\sigma]]_F
 =\begin{cases}
 [S_{1,1,F}],&\text{if $b$ has defect zero},\\
 0,&\text{if $b$ has positive defect}
 \end{cases}
 \quad\text{in }K_0(\mathcal F\mathcal B\ell_{F,k}).
\]
\end{corollary}

\begin{proof}
Boltje--Bouc--Y\i lmaz first prove that, for every group--block pair,
the alternating sum above is an integer multiple of $[S_{1,1,F}]$:
every other simple-functor multiplicity cancels
\cite[Theorem~3.5\textup{(1)}]{BBY2026}. The universal numerical
blockwise Alperin weight conjecture determines the remaining coefficient
through its chain formulation. Their Theorem~3.5\textup{(2)} therefore
gives the stated identity from \Cref{thm:main}.
\end{proof}
